\begin{tikzpicture}[
    >=Latex,
    font=\small,
    node distance=8mm and 12mm
]

% ============================================================
% PANEL (a)
% ============================================================

\node[
    anchor=west,
    font=\bfseries
] (titlea) at (0,0)
    {(a) Internet routing: geographic distance vs.\ routing path};

% Endpoints
\node[
    endpoint,
    below=7mm of titlea,
    xshift=-65mm
] (sourceA)
    {Source\\(host / AS X)};

\node[
    endpoint,
    right=110mm of sourceA
] (destA)
    {Destination\\(host / AS Y)};

% Autonomous systems
\node[
    autonomoussystem,
    right=10mm of sourceA
] (as1) {AS 1};

\node[
    autonomoussystem,
    right=17mm of as1,
    yshift=10mm
] (as2) {AS 2};

\node[
    autonomoussystem,
    right=17mm of as2
] (as4) {AS 4};

\node[
    autonomoussystem,
    right=17mm of as4,
    yshift=-10mm
] (as5) {AS 5};

% Transit AS
\node[
    autonomoussystem,
    below=15mm of as4
] (as6) {AS 6\\(transit)};

% IXP
\node[
    ixp,
    right=8mm of as2
] (ixp) {IXP};

% Links
\draw[link] (sourceA) -- (as1);
\draw[link] (as1) -- (as2);
\draw[link] (as2) -- (ixp);
\draw[link] (ixp) -- (as4);
\draw[link] (as4) -- (as5);
\draw[link] (as5) -- (destA);

% Transit path
\draw[transit] (as1) -- (as6);
\draw[transit] (as6) -- (as5);

% Actual BGP-selected path
\draw[bgppath] (as1) -- (as2);
\draw[bgppath] (as2) -- (ixp);
\draw[bgppath] (ixp) -- (as4);
\draw[bgppath] (as4) -- (as5);

% Geographic shortest path
\draw[shortest]
    (sourceA.north east)
    .. controls
        +(25mm,12mm)
        and +(-25mm,12mm)
    ..
    (destA.north west);

\node[
    text=green!45!black,
    font=\footnotesize\bfseries,
    align=center
] at (
    $(sourceA.north)!0.5!(destA.north)+(0,15mm)$
)
{
    Geographically shortest path\\
    (not necessarily used)
};

\node[
    text=blue!80!black,
    font=\footnotesize\bfseries,
    align=center
] at (
    $(as2.south)!0.5!(as4.south)+(0,-9mm)$
)
{
    Actual routing path selected by BGP\\
    based on policies, peering, and traffic engineering
};

% ============================================================
% Key concepts box
% ============================================================

\node[
    draw,
    rounded corners=3pt,
    fill=blue!3,
    text width=61mm,
    align=left,
    inner sep=5pt,
    font=\footnotesize
] at (
    $(sourceA.south)!0.5!(as6.south)+(-2mm,-17mm)$
)
{
    \textbf{Key concepts}
    \begin{itemize}
        \setlength{\itemsep}{1pt}
        \setlength{\parskip}{0pt}
        \item Internet routing (e.g., BGP) selects paths according
        to routing policies, not physical distance.
        
        \item Peering relationships and traffic engineering may
        lead to longer geographic paths.
        
        \item Network proximity does not necessarily imply
        geographic proximity.
    \end{itemize}
};

% Separator
\draw[
    gray!50,
    line width=0.7pt
] (-88mm,-67mm) -- (88mm,-67mm);

% ============================================================
% PANEL (b)
% ============================================================

\node[
    anchor=west,
    font=\bfseries
] (titleb) at (-88mm,-73mm)
{
    (b) Anycast routing: one IP, multiple geographically
    distributed replicas
};

% Source
\node[
    endpoint,
    below=7mm of titleb,
    xshift=-65mm
] (sourceB)
    {Source\\(host / AS X)};

% ASes
\node[
    autonomoussystem,
    right=10mm of sourceB
] (bas1) {AS 1};

\node[
    autonomoussystem,
    right=17mm of bas1,
    yshift=10mm
] (bas2) {AS 2};

\node[
    autonomoussystem,
    right=17mm of bas2
] (bas4) {AS 4};

\node[
    autonomoussystem,
    right=17mm of bas4,
    yshift=-10mm
] (bas5) {AS 5};

% IXP
\node[
    ixp,
    right=8mm of bas2
] (bixp) {IXP};

% Anycast IP
\node[
    anycastip,
    right=12mm of bas5
] (ip)
{
    Anycast IP\\
    \texttt{203.0.113.1}
};

% Replicas
\node[
    replica,
    right=16mm of ip,
    yshift=13mm
] (r1)
{
    Replica 1\\
    Madrid, Spain
};

\node[
    replica,
    right=16mm of ip
] (r2)
{
    Replica 2\\
    Paris, France
};

\node[
    replica,
    right=16mm of ip,
    yshift=-13mm
] (r3)
{
    Replica 3\\
    New York, USA
};

% Links
\draw[link] (sourceB) -- (bas1);
\draw[link] (bas1) -- (bas2);
\draw[link] (bas2) -- (bixp);
\draw[link] (bixp) -- (bas4);
\draw[link] (bas4) -- (bas5);
\draw[link] (bas5) -- (ip);

% Selected routing path
\draw[bgppath]
    (bas1) -- (bas2)
    -- (bixp)
    -- (bas4)
    -- (bas5);

% Anycast replicas
\draw[replicaLink] (ip) -- (r1);
\draw[replicaLink] (ip) -- (r2);
\draw[replicaLink] (ip) -- (r3);

% Replica annotation
\node[
    align=left,
    font=\footnotesize
] at (
    $(r2.east)+(20mm,0)$
)
{
    Multiple replicas\\
    sharing the same\\
    anycast IP address
};

% ============================================================
% Anycast principle box
% ============================================================

\node[
    draw,
    rounded corners=3pt,
    fill=violet!3,
    text width=68mm,
    align=left,
    inner sep=5pt,
    font=\footnotesize
] at (
    $(sourceB.south)!0.5!(bas5.south)+(-2mm,-18mm)$
)
{
    \textbf{Anycast principle}
    \begin{itemize}
        \setlength{\itemsep}{1pt}
        \setlength{\parskip}{0pt}
        \item Multiple geographically distributed replicas share
        the same IP address.
        
        \item Routing protocols deliver packets to the replica
        selected according to routing policies and current
        network conditions.
        
        \item The selected replica, and therefore the geographic
        destination, may vary across sources and over time.
    \end{itemize}
};

\end{tikzpicture}
